\subsection{Further attempts at the two remaining torus questions}

\statement{C.24}{Proved}{The proposed cubic lattice admits an invariant unimodular surface form}
Let $q\equiv1\pmod3$ be an odd prime.  The
$\mathbb Z_3[C_3]$-module
\[
 L_q=\mathbb Z_3[\zeta_3]\oplus
                       \mathbb Z_3[C_3]^{\,q-1}             \tag{C.24.1}
\]
admits a $C_3$-invariant unimodular alternating
form.  All integral fixed vectors have pairings
divisible by three.  This module is realized
by the homology of a geometric order-three
action on a closed surface of genus
$(3q-1)/2$.  Thus a contradiction based solely
on the module (C.24.1), its surface form and
mod-three isotropy of integral fixed characters
cannot exclude a scalar cubic torus coordinate.
This is an algebraic limitation of that
approach, not a realization of such a torus
coordinate.

\begin{proof}
On the cyclotomic summand take
\[
 T=\begin{pmatrix}0&-1\\1&-1\end{pmatrix},
 \qquad J=\begin{pmatrix}0&1\\-1&0\end{pmatrix}.
\]
Then $T^3=I$, $T^2+T+I=0$, and
$T^{\mathsf T}JT=J$.  This summand has no
nonzero integral fixed vector.
Since $q-1$ is even, pair the regular modules
in $(q-1)/2$ pairs.  On a pair with bases
$e_0,e_1,e_2$ and $f_0,f_1,f_2$, let the
generator cyclically permute each basis and
put $\langle e_i,f_j\rangle=\delta_{ij}$,
with the reverse pairing negative and all
other pairings zero.  This form is invariant
and unimodular.  The fixed submodule is
generated by $e_0+e_1+e_2$ and
$f_0+f_1+f_2$; their pairing is three.
Orthogonal direct sum proves both lattice
assertions.

For the surface realization, begin with an
elliptic curve admitting multiplication by
$\zeta_3$, whose order-three action has three
fixed points and cyclotomic integral homology.
Choose a free orbit of three small discs and
attach three handles equivariantly.  This
adds one pair of regular homology modules:
the three handle $a$ cycles and the three
handle $b$ cycles are cyclically permuted and
have the displayed hyperbolic pairings.
Repeat $(q-1)/2$ times, using disjoint free
orbits.  Genus becomes
$1+3(q-1)/2=(3q-1)/2$, and homology is
(C.24.1).  The three original fixed points
remain.  This construction has global
orientation multiplier one and does not
descend to the hypothetical scalar torus
action.  It proves exactly the asserted
compatibility of the linear data.
\end{proof}


\statement{C.25}{Proved}{A free-factor route would settle the actual fixed kernel}
In the notation of C.12, if the actual other
fixed group $B$ is a free profinite factor
of $F=\widehat F_2$, then $B\cong\widehat{\mathbb Z}$
and $K=\ker(\chi|_B)=1$.  The abstract
projective retract constructed in C.2 is
not a free profinite factor of $\widehat F_2$.

\begin{proof}
Suppose $F=B\amalg D$.  Both factors are
retracts of the finitely generated group
$F$ and are therefore finitely generated.
On abelianization this is the direct sum
$F^{\rm ab}=B^{\rm ab}\oplus D^{\rm ab}$.
The inclusion of $B^{\rm ab}$ is injective
and lies in the fixed module of the
negative involution, since $B$ is its
actual fixed subgroup.  At each prime this
fixed module has rank one: the order-two
matrix has determinant $-1$ and eigenvalues
$1,-1$.  The epimorphism
$\chi:B\to\widehat{\mathbb Z}$ forces
$B^{\rm ab}$ to have rank at least one
at every prime.  As a direct summand of
$\widehat{\mathbb Z}^2$, it is consequently
$\widehat{\mathbb Z}$, and so is
$D^{\rm ab}$.

For every finite group $A$, the universal
property of the full profinite free product
gives the exact count
\[
 |\operatorname{Hom}(B,A)|
 |\operatorname{Hom}(D,A)|=|A|^2.                         \tag{C.25.1}
\]
Each factor on the left is at least $|A|$,
because homomorphisms through its abelianization
$\widehat{\mathbb Z}$ are parametrized by the
image of one topological generator.
A homomorphism of $B$ with nonabelian image
would add another homomorphism beyond these
$|A|$, contradicting (C.25.1).  Hence every
finite image of $B$ is abelian, so $B$ is
abelian and C.1 gives $K=1$.

For the C.2 model the abelianization is also
$\widehat{\mathbb Z}$, but it has twelve
homomorphisms to $S_3$.  To see this, its
pro-three generators $x,y$ map into $A_3$.
The prime-to-three generator $z$ maps to
identity or a transposition.  If it is
identity, $zyz^{-1}=y^{-1}$ forces $y=1$
and gives three choices for $x$.
For each of its three transposition values,
$zxz^{-1}=x$ forces $x=1$ and gives three
choices for $y$.  Thus there are
$3+3\cdot3=12$ maps.  A free complement
would have procyclic abelianization and
give at least six maps to $S_3$, contradicting
$12\cdot6>6^2$.  The finite enumeration in
\path{code/check_cubic_and_factor.py} verifies
this count and the lattice matrices in C.24.
\end{proof}

\statement{C.26}{Open}{The scalar cubic existence question in prompt C.3}
We have not constructed or excluded an
order-three element of the full torus
loop-outer group whose multiplier has a
nonidentity cubic coordinate.  C.16 is
specific to the three-cusp structure and
does not prove the torus assertion.
C.24 explains exactly why replacing the
quartic involution in C.22 by the proposed
cubic integral module does not yield the
same contradiction: the fixed vectors
come from norms in regular modules and
their pairings vanish modulo three.
The module is even geometrically realizable
on the filled covering surface, while the
needed descent to the original punctured
torus with its specified nontrivial multiplier
is absent.

An additional cyclic-cover attempt has the
following precise limit.  At a scalar cubic
coordinate $q$, an eigencharacter
$F\to C_q$ is preserved up to a primitive
cube-root scalar, rather than pointwise.
The induced order-three action consequently
permutes nontrivial deck eigenspaces in
orbits of length three.  Its commutant is
no longer the product of unramified fields
with multiplicity one used in C.16; the
resulting permutation blocks themselves
have order three.  Thus absence of a
nontrivial cubic root in an unramified
triadic field does not exclude this different
semilinear action.
The finite Frobenius quotients proposed in
LM4 D.71 and D.73 would impose stronger
full-group constraints; without a separate
proof that those constraints are incompatible,
or a construction of a compatible full
action, they do not answer this row.

\statement{C.27}{Proved}{Mixed-sign diagonal power involutions fail the full torus cusp test}
For $u\in\widehat{\mathbb Z}^{\times}$ with $u^2=1$, the
continuous automorphism
\[
 \alpha_u:\widehat F(x,y)\longrightarrow\widehat F(x,y),
 \qquad x\longmapsto x^u,\quad y\longmapsto y^{-u}           \tag{C.27.1}
\]
has order two and determinant $-1$ on full abelianization.
It preserves the torus cusp $[x,y]$ up to conjugacy and
the power $-1$ if and only if $u=1$ or $u=-1$ globally.
In these two cases it is a standard reflection and its
noncusp fixed group is $\widehat{\mathbb Z}$.

Every pro-$p$ image of (C.27.1) is a standard reflection,
including when the full cusp assertion fails.  Thus this
family gives an explicit mixed-prime obstruction to
one natural attempted construction in prompt C.2.
It does not classify general negative torus involutions.

\begin{proof}
Each primary coordinate $u_p$ equals $1$ or $-1$, since
$\mathbb Z_p$ is an integral domain.  The universal
property of the full free group defines (C.27.1), and
applying it twice fixes both generators.  Its determinant
is $-u^2=-1$, so it is not the identity.  At any one prime
it fixes one generator and inverts the other.  This
proves the assertions about the actual automorphism and
its pro-$p$ images.

Suppose first that the signs at two odd primes differ.
Name them $q,p$ so that $u_q=1$ and $u_p=-1$.  Put
$d=\operatorname{ord}_p(q)$, and choose an element
$\zeta$ of order $p$ in $\mathbb F_{q^d}^{\times}$.
Its degree over $\mathbb F_q$ is $d$.  Form the finite
group
\[
 G=(\mathbb F_{q^d},+)\rtimes C_p,
 \qquad (v,i)(w,j)=(v+\zeta^i w,i+j).
                                                               \tag{C.27.2}
\]
The elements $X=(1,0)$ and $Y=(0,1)$ generate $G$:
the conjugates of $X$ span the field additively because
$\zeta$ has degree $d$.  They have orders $q,p$,
respectively.  The epimorphism $x\mapsto X,y\mapsto Y$
therefore sends both $\alpha_u(x)$ and $\alpha_u(y)$
to the original generator images.  However
\[
 [X,Y]=(1-\zeta,0)
\]
is not conjugate to its inverse.  Conjugation on this
nonzero additive element multiplies it by a power of
$\zeta$, whereas $-1$ is not a power of an element of
odd order.  Thus the required full cusp conjugacy is
impossible in this finite quotient.

It remains to exclude a common sign at all odd primes
with the opposite sign at two.  In $\operatorname{SL}_2
(\mathbb F_7)$ take
\[
 X=\begin{pmatrix}0&6\\1&3\end{pmatrix},\qquad
 Y=\begin{pmatrix}0&2\\3&6\end{pmatrix},\qquad
 C=[X,Y]=\begin{pmatrix}5&6\\1&0\end{pmatrix}.              \tag{C.27.3}
\]
They have orders $8,3,14$, respectively.  In particular
$X^4=C^7=-I$.  The element $-C=I+N$ has $N\ne0$,
$N^2=0$.  Conjugacy of $C$ to $C^{-1}$ would conjugate
$N$ to $-N$.  In a basis adapted to the kernel of
$N$, a determinant-one conjugator has diagonal entries
$a,a^{-1}$, and changes the nonzero upper-right entry
of $N$ by $a^2$.  It would follow that $a^2=-1$ in
$\mathbb F_7$, which is false.  Hence $C$ is not
conjugate to $C^{-1}$ in this finite group.

If $u_2=1$ and all odd signs are $-1$, (C.27.1) fixes
both matrices in (C.27.3).  If $u_2=-1$ and all odd
signs are $1$, it inverts both.  In the latter case
\[
 [X^{-1},Y^{-1}]=(XY)^{-1}[X,Y](XY),
\]
so the commutator image still lies in the conjugacy
class of $C$, not that of its inverse.  Both remaining
mixed patterns are therefore excluded.

For $u=1$, the identity
$[x,y^{-1}]=y^{-1}[x,y]^{-1}y$ proves cusp preservation;
for $u=-1$ use
$[x^{-1},y]=x^{-1}[x,y]^{-1}x$.
These reflections are conjugate by the basis exchange.
For the first one, with section generator $s$, setting
$t=ys$ gives the full presentation
\[
 \widehat F_2\rtimes_{\alpha_1} C_2
       =(\widehat{\mathbb Z}\times C_2)\amalg C_2,
 \qquad [x,s]=1,\quad s^2=t^2=1.
\]
The centralizer theorem for profinite free products,
\cite[Theorem B']{A_HR1985}, identifies its noncusp
geometric centralizer with $\overline{\langle x\rangle}$.
The other case exchanges $x,y$.  Thus $K'=1$ for
this family whenever it is actually peripheral.

The executable \path{code/check_diagonal_involutions.py}
checks the order-$21$ instance of (C.27.2), enumerates
all $336$ matrices of the group in (C.27.3), verifies
that $X,Y$ generate it, and checks the two disjoint
commutator conjugacy classes.  The general argument
above supplies the assertion for every sign pattern.
\end{proof}

\statement{C.28}{Proved}{Self-centralizing sections for every pure negative sphere action}
Let $n\geq2$, and let
\[
 F_n=\widehat{\pi_1(S_{0,n+1})}
   =\widehat F(x_1,\ldots,x_n),\qquad
 x_{n+1}=(x_1\cdots x_n)^{-1}.
\]
Suppose an outer involution preserves each of the $n+1$
peripheral conjugacy classes and sends each $x_j$ to a
conjugate of $x_j^{-1}$.  In its centre-free extension
$1\to F_n\to E\to C_2\to1$, there are exactly $n+1$
$F_n$-conjugacy classes of sections.  Every section
$S$ satisfies
\[
                         C_{F_n}(S)=1.                 \tag{C.28.1}
\]
In particular every genuine involutory representative has
trivial fixed subgroup.  The assertion concerns the full
profinite group, not only its pro-$2$ quotient.

More generally, let an involution act compatibly on the full
profinite pure mapping class groups of the punctured spheres
in a forgetful tower, preserving the individual peripheral
classes of each point-pushing kernel with exponent $-1$.
Its fixed subgroup in every group of this tower is trivial.
The same holds in the pure braid tower, provided the
two-strand group is inverted.  In particular, the actual
action of any multiplier-$-1$ involution of
$\widehat{GT}$ on a pure braid group has trivial fixed
subgroup.

\begin{proof}
Each peripheral generator has primitive homology.  Thus
C.5 identifies its geometric normalizer with
$L_j=\overline{\langle x_j\rangle}\cong\widehat{\mathbb Z}$.
A lift in $N_E(L_j)$ acts on $L_j$ by inversion; its square
belongs to $L_j$ and is fixed by inversion, and hence is one.
Consequently the extension splits and
\[
               D_j=N_E(L_j)=L_j\rtimes_{-1}C_2.
\]
Fix a splitting $E=F_n\rtimes\langle v\rangle$ and put
$k=\mathbb F_2$.  Let $u:E\to k$ be its quotient character.
The action on $H_1(F_n,\mathbb Z_2)$ is $-I$.  For the
integral coordinate character $\varphi_i:F_n\to\mathbb Z_2$,
its reduction extends to $b_i:E\to k$ with $b_i(v)=0$.
The cochain
$\widetilde b_i(hv^\epsilon)=\varphi_i(h)\pmod4$
has coboundary $2\epsilon\varphi_i(h')$ on
$(hv^\epsilon,h'v^\eta)$.  The mod-two Bockstein therefore
gives
\[
                            b_i^2=ub_i.                 \tag{C.28.2}
\]
The Hochschild--Serre spectral sequence has just two fibre
rows.  The section makes inflation injective, so its
transgression vanishes; multiplication by $u$ makes all
positive-column differentials vanish.  Thus, in every
sufficiently high degree, cohomology has basis
$u^d,u^{d-1}b_1,\ldots,u^{d-1}b_n$.  The degree-zero
localized algebra consequently has basis
$1,B_1,\ldots,B_n$, where $B_i=b_i/u$.  Every element is
idempotent by (C.28.2), so this finite Boolean algebra is
\[
                     H^\bullet(E,k)[u^{-1}]_0
                           \cong k^{\,n+1}.              \tag{C.28.3}
\]

Apply the full finite-coefficient restriction theorem
used in C.7--C.8, namely Theorem~2 of
\cite{CEKochloukovaZalesskii2007}.  Its hypotheses hold:
the open geometric group has cohomological dimension one,
and every finite subgroup injects into $C_2$.  The dense
restriction map is an isomorphism in high degrees because
its source is finite.  For trivial coefficients its
localized factors are
\[
       k\oplus H^1(C_{F_n}(S),k)u^{-1},
\]
whose second summands have square zero.  Comparing with
the reduced algebra (C.28.3) gives exactly $n+1$ section
classes, with zero first mod-two cohomology for each
geometric centralizer.  The $n$ values $B_i(S)$ distinguish
these classes, since $1,B_1,\ldots,B_n$ span (C.28.3).
This scalar calculation alone would not exclude a
nontrivial pro-odd centralizer.

The needed full-coefficient information comes from the
actual peripheral incidence.  Form a finite graph whose
vertices are the $n+1$ global section classes and whose
$j$th edge joins the images of the two local section
classes of $D_j$.  Their labels in $k^n$ differ by
$e_j$ for $j\leq n$, and by $e_1+\cdots+e_n$ for
$j=n+1$.  Indeed the local representatives differ by
the corresponding $x_j$, and evaluation of the $b_i$
gives exactly its homology class.  None of these
differences is zero.  The only nonempty subset of the
$n+1$ edge directions whose sum is zero is the whole
set.  A graph with $n+1$ edges and $n+1$ vertices
has a cycle; the directions around any cycle sum to
zero.  That cycle must therefore use every edge and
every vertex.  Each global section class occurs at
exactly two peripheral endpoints.

Now apply C.6 to the sphere pair and an arbitrary finite
self-dual permutation $E$-module $M$.  Brown restriction
decomposes the localized absolute group into the section
normalizers $C_{F_n}(S)\times S$.  Boundary restriction
is the sum, for each section class, of its restriction
to $S$ twice, with the actual conjugating coefficient
transports.  C.6 says that the total rank is half the
boundary dimension.  Since each individual restriction
rank is at most its target dimension, and each occurs
twice, every such individual restriction is onto.

If $1\ne h\in C_{F_n}(S)$, choose a finite quotient
$D$ of $E$ in which the image of $h$ is outside the
image $\overline S$ of $S$.  This is possible since
$h\notin S$ and $S$ is finite and closed.  Take
$M=k[D/\overline S]$.  Positive cyclic cohomology
has the $S$-fixed cosets as a basis.  The element $h$
moves the identity-coset basis vector to a distinct
fixed-coset basis vector.  The image of restriction
from $C_{F_n}(S)\times S$, however, is fixed by this
centralizer, by naturality under inner conjugation.
This contradicts the established surjectivity, proving
(C.28.1).  No presentation of $E$ as a profinite
free product has been used.

Finally the full profinite Birman sequence for ordered
punctures is
\[
 1\longrightarrow\widehat{\pi_1(S_{0,m-1})}
 \longrightarrow\widehat\Gamma_{0,m}
 \longrightarrow\widehat\Gamma_{0,m-1}
 \longrightarrow1,\qquad m\geq4 .
\]
Its full exactness and topology comparison are the
genus-zero configuration comparison recalled in B.30,
and treated in \cite{CEBoggiCongruence2020}.

For a compatible involution satisfying the stated
peripheral hypothesis, a fixed element maps to a
fixed element in the quotient.  Induction from
$\Gamma_{0,3}=1$ puts it in the point-pushing kernel,
where (C.28.1) makes it one.  This last assertion
requires an actual compatible action on the whole
tower; no action of $\widehat{GT}_0$ on the
five-point group has been presumed.

For the pure braid version use the actual full sequences
\[
 1\longrightarrow\widehat F_{m-1}
 \longrightarrow\widehat P_m
 \longrightarrow\widehat P_{m-1}
 \longrightarrow1,\qquad m\geq3,
\]
starting from $\widehat P_2=\widehat{\mathbb Z}$.
Here $P_m$ denotes the discrete pure braid group on
$m$ ordered strands.  To verify full exactness directly,
the discrete Fadell--Neuwirth sequence splits.  Given a
characteristic finite-index subgroup $V$ of its finitely
generated free kernel, the quotient action on $F_{m-1}/V$
has finite image.  If $Q_0$ is the kernel of that action,
$V\rtimes Q_0$ is finite-index in
$F_{m-1}\rtimes P_{m-1}$ and meets the free kernel in $V$.
Such $V$ are cofinal, proving its full induced topology;
the split quotient then gives the displayed exact sequence
after completion.  No pro-$p$ replacement is involved.

For the actual operadic action of $\widehat{GT}$
\cite[Sections 7 and 11]{CEBHR2019}, finite
point-pushing cusps are sent to conjugates of their
multiplier powers.  The infinity cusp is
$z_m z_{m-1}^{-1}$, where $z_j$ is the full twist on
the first $j$ strands.  Choose the parenthesization that
groups the first $m-1$ strands.  Full twists are powered
by the multiplier, including the embedded $z_{m-1}$;
since $z_m$ is central, their displayed quotient is
inverted when the multiplier is $-1$.  These are the
actual individual peripheral classes needed for the
free-kernel assertion.  Transport to another
parenthesization only conjugates these classes.
Moreover the pure quotient acts on the point-pushing
kernel preserving each of its peripheral conjugacy
classes, so conjugacy in the ambient pure braid group
here is already conjugacy inside that kernel.
The base
$\widehat P_2$ has no nonidentity element fixed by
inversion.  Induction proves the pure-braid assertion,
in particular for $\widehat P_4$.
\end{proof}

\statement{C.29}{Proved}{A full boundary test for automorphisms}
Let $S$ be a connected compact oriented surface with $m\geq1$ boundary components and negative Euler characteristic. Let $F=\widehat{\pi_1(S)}$ be its full profinite fundamental group and let $c_1,\ldots,c_m$ be its positively oriented boundary generators, based by chosen arcs. Suppose a continuous endomorphism $\varphi:F\to F$ satisfies
\[
 \varphi(c_j)=g_jc_{\sigma(j)}^{\mu_j}g_j^{-1}
 \quad(1\leq j\leq m),
 \qquad \sigma\in\mathfrak S_m,\quad
 \mu_j\in\widehat{\mathbb Z}^{\times},\quad g_j\in F.
\]
Then $\varphi$ is an automorphism. If $m\geq2$, its exponents $\mu_j$ are necessarily equal. In particular, for the full profinite pants group, a substitution
\[
 x\longmapsto x^\mu,\qquad y\longmapsto u^{-1}y^\mu u
\]
with $\mu\in\widehat{\mathbb Z}^{\times}$ is an automorphism as soon as $x^\mu u^{-1}y^\mu u$ is conjugate to $(xy)^\mu$. No separate full-generation hypothesis is needed.

\begin{proof}
Suppose the compact image $H=\varphi(F)$ is proper. Choose a proper open subgroup $U\leq F$ containing it. Writing $G=\pi_1(S)$, the subgroup $U_0=U\cap G$ has finite index $d=[F:U]>1$, and $\widehat{U_0}=U$. Let $S_U\to S$ be its connected finite covering. These are full finite-index topology comparisons, not pro-$p$ comparisons.

Replace each profinite conjugator by a discrete coset representative. Density of $G$ and openness of $U$ give $a_j\in G\cap Ug_j$. Write $a_j=v_jg_j$ with $v_j\in U$. The given image of $c_j$ is then $U$-conjugate to $a_j c_{\sigma(j)}^{\mu_j}a_j^{-1}$. As this lies in $U$ and $\mu_j$ is a full unit, closedness gives $a_j c_{\sigma(j)}a_j^{-1}\in U$. This discrete element belongs to $U_0$. The corresponding boundary component $b_j$ of $S_U$ has covering degree one over the $\sigma(j)$th base boundary, and in $H_1(U,\widehat{\mathbb Z})$ we have
\[
                       [\varphi(c_j)]=\mu_j[b_j].
\]
The selected $b_j$ are distinct because their base boundaries are distinct.

The oriented surface relation is a product of commutators times $c_1\cdots c_m=1$. Its image in the abelianization of $U$ therefore gives
\[
                         \sum_{j=1}^m\mu_j[b_j]=0.
\]
If $S_U$ has $r$ boundary components, their classes in its first integral homology have precisely one relation: their total oriented sum is zero. Equivalently, the kernel of the boundary-class map $\mathbb Z^r\to H_1(S_U,\mathbb Z)$ is generated by $(1,\ldots,1)$. This remains true over $\widehat{\mathbb Z}$: choose the usual handle and all-but-one boundary basis, so the statement follows by an explicit split free-module calculation, without an exactness assumption on an arbitrary completion. Consequently the coefficient vector having entries $\mu_j$ on the selected components and zero on every unselected component must be constant. If there were any unselected component, the constant would be zero, contradicting that every $\mu_j$ is a unit. Thus the $b_j$ are all boundary components of $S_U$.

There is therefore exactly one boundary component over each boundary of $S$, and its covering degree is one. The sum of boundary covering degrees over any one base boundary equals $d$, by the ordinary finite covering-space formula. Hence $d=1$, a contradiction. Thus $\varphi$ is surjective. A finitely generated profinite group is Hopfian, so it is an automorphism. Alternatively, preimage under a surjection permutes the finite set of open subgroups of each bounded index; every element of its kernel then lies in every open subgroup and is one.

For $m\geq2$, apply the same unique boundary relation directly in $H_1(F,\widehat{\mathbb Z})$: it gives $\mu_1=\cdots=\mu_m$. The pants specialization has $c_3=(xy)^{-1}$, and its stated third-cusp equation supplies exactly the final boundary condition. This is an independent covering proof of the pants generation step already asserted in LM4 D.35, which used the closed-subgroup rank test of LM4 D.32. It also generalizes the direct punctured-torus boundary test independently proved in C.13. No new resolution of the existence of a third-cusp parameter is claimed.
\end{proof}


\statement{C.30}{Proved}{Unit hexagon solutions are full automorphisms}
Let $F=\widehat F(x,y)$, put $z=(xy)^{-1}$, and let
$\lambda\in\widehat{\mathbb Z}^{\times}$ and $f\in F'$.
Write $m=(\lambda-1)/2\in\widehat{\mathbb Z}$.
Suppose the following are identities in the full profinite group:
\[
 \begin{split}
 f(x,y)f(y,x)&=1,\\
 f(x,y)x^m f(z,x)z^m f(y,z)y^m&=1.
 \end{split}                                                    \tag{C.30.1}
\]
Then the continuous endomorphism
\[
             \varphi(x)=x^\lambda,\qquad
             \varphi(y)=f^{-1}y^\lambda f
\]
is an automorphism. In particular no separate invertibility hypothesis
is needed for a unit pair satisfying these two full equations. If it
also satisfies the full pentagon, it belongs to the full
Grothendieck--Teichm\"uller group, including its unit requirement.

\begin{proof}
This is an algebraic consequence of C.29, but it is essential not to
assume invertibility while deriving the third-cusp equation.
Set $A=f(x,y)$, $B=f(y,z)$ and $C=f(z,x)$. The second equation in
(C.30.1) first gives
\[
                C z^m B=x^{-m}A^{-1}y^{-m}.                      \tag{C.30.2}
\]
Apply that same identity under the continuous substitution
$(x,y)\mapsto(x,z)$. Its third argument is
\[
                w=(xz)^{-1}=xyx^{-1}=z^{-1}yz.
\]
The first equation in (C.30.1), applied under substitutions, and
conjugation-equivariance of profinite words give
\[
 f(x,z)=C^{-1},\qquad
 f(w,x)=xA^{-1}x^{-1},\qquad
 f(z,w)=z^{-1}B^{-1}z.
\]
Consequently the substituted hexagon becomes
\[
 C^{-1}x^{m+1}A^{-1}y^{m+1}B^{-1}z^{m+1}=1,
\]
where $x^{-1}z^{-1}=y$ was used. Rearrangement gives the second identity
\[
          C z^{-m-1}B=x^{m+1}A^{-1}y^{m+1}.                     \tag{C.30.3}
\]
Multiply (C.30.3) by the inverse of (C.30.2). Since $2m+1=\lambda$,
then conjugate the resulting equality by $x^m$. One obtains
\[
 x^\lambda A^{-1}y^\lambda A
       =(x^mC)(xy)^\lambda(x^mC)^{-1}.                         \tag{C.30.4}
\]
All substitutions used here are legitimate continuous homomorphisms
out of the free profinite group; none has been assumed invertible.
Equation (C.30.4) is the third peripheral condition, and C.29 makes
$\varphi$ a full automorphism.

For completeness the pentagon assertion concerns actual units, not
only endomorphisms. The normalized two-equation pairs are now
contained in $\operatorname{Aut}(F)$. Their pentagon solutions form
a closed submonoid there, by the parenthesized-braid presentation
\cite[Lemma 7.1 and Proposition 7.3]{CEBHR2019}. The automorphism group
of a finitely generated profinite group is profinite: characteristic
open subgroups are cofinal, and their finite automorphism groups
give its topology. Thus any closed submonoid is a group, since
$\varphi^{n!-1}\to\varphi^{-1}$. This supplies the inverse within
the same pentagon solution set.
\end{proof}

\paragraph{Source discrepancy and scope.}
Lochak and Schneps explicitly asked in 2006 whether a unit normalized
pair is automatically invertible, and then asked the same question
with their relations (I), (II), (III) imposed
\cite[Question 1.1]{A_LS2006}. Fresse's later survey also distinguishes
the profinite invertibility issue from the Malcev case
\cite[\S1.1, manuscript p.~14]{CEFresse2019}.
C.30 answers the constrained version by the displayed full-group
calculation and C.29's finite-cover proof. It does not use the
pro-$p$ Frattini criterion, infer invertibility merely from
abelianization, or claim a literature priority. It also does not
prove that a bare profinite pentagon has a unit parameter or implies
the hexagon; those are different requirements in A.37--A.39.

The unconstrained version of that old question has a small finite
counterexample, which also calibrates this source comparison.
Take $\lambda=1$, $f=[x,y]^2$ and the quotient onto $A_5$ with
$x=(123)$ and $y=(345)$, using $[x,y]=xyx^{-1}y^{-1}$ and
right-to-left permutation products. Then $f$ has image $(134)$,
and $f^{-1}yf$ has image $(135)$. The original pair generates
$A_5$, whereas the substituted pair generates the order-twelve
point stabilizer fixing $4$. Thus this unrestricted normalized
endomorphism is not surjective. The exact finite calculation is
recorded in \path{code/check_unconstrained_substitution.py}.

\statement{C.31}{Proved}{Every full braid-operad endomorphism has a unit multiplier}
Let $\Phi$ be an object-fixing continuous endomorphism of the symmetric
operad $\widehat{\mathrm{PaB}}$ of full profinite parenthesized braids.
Let $\lambda\in1+2\widehat{\mathbb Z}$ be its exponent on the elementary
two-strand braiding. Then
\[
                         \lambda\in\widehat{\mathbb Z}^{\times}.
\]
Consequently the full Grothendieck--Teichm\"uller endomorphism monoid,
defined by the full duality, hexagon and pentagon equations with an odd
parameter, is already its group of units.

\begin{proof}
We first establish the precise full-completion input. Write $B_n$ and
$P_n$ for the ordinary and pure discrete Artin braid groups. The
finite-index inclusion $P_n\leq B_n$ induces its full profinite
topology, by taking cores of finite-index subgroups in $B_n$. Hence
\[
 1\longrightarrow\widehat P_n\longrightarrow\widehat B_n
       \longrightarrow S_n\longrightarrow1                     \tag{C.31.1}
\]
is exact with the full completion in the first term. Likewise the
split Fadell--Neuwirth sequence
$1\to F_{n-1}\to P_n\to P_{n-1}\to1$ retains the full topology on
its kernel. Indeed any finite-index subgroup of the finitely
generated free kernel contains a characteristic finite-index subgroup
$W$; the action on $F_{n-1}/W$ factors through a finite quotient of
$P_{n-1}$ and gives a finite semidirect-product quotient detecting
that kernel quotient.

Finite-rank free groups are good. Comparison of the discrete and
continuous Hochschild--Serre spectral sequences therefore gives
goodness of $P_n$ inductively, and then of $B_n$ using (C.31.1).
The finite-coefficient cohomology of the kernels in this induction is
finite, so the comparison applies to every required coefficient
module. These are also the precise primary-source conclusions of
\cite[Proposition 5.10 and Corollary 5.11]{CEHorel2017} and
\cite[\S2.3, Proposition 1]{CEMarin2012}.
The unordered configuration space of $n$ distinct points in
$\mathbb C$ is an aspherical smooth manifold of real dimension $2n$:
asphericity follows from the ordered Fadell--Neuwirth fibrations,
and passage to the free finite symmetric-group quotient preserves
the universal cover. Its CW dimension is at most $2n$, so
$\operatorname{cd}(B_n)\leq2n$. Goodness now implies, for every
finite continuous $p$-primary module $A$,
\[
 H^j_{\mathrm{cts}}(\widehat B_n,A)=0\qquad(j>2n).
\]
An arbitrary discrete $p$-primary module is the filtered union of
its finite invariant submodules, and continuous cochains on compact
profinite powers have finite image. Thus the same vanishing holds
for every such module and
$\operatorname{cd}_p(\widehat B_n)\leq2n$ for every prime $p$.
A closed $C_p$ is impossible, since cohomological dimension does
not increase on closed subgroups and
$H^j(C_p,\mathbb F_p)\ne0$ in every degree. In particular
\[
                     \widehat B_n\text{ is torsion-free}.        \tag{C.31.2}
\]
This implication is also recorded in
\cite[Lemma 8.4]{CEMorales2024}. It uses goodness and actual finite
cohomological dimension, rather than virtual cohomological dimension
or torsion-freeness of the discrete group alone.

Choose a binary parenthesization $b$ in arity $n$. The full subgroupoid
on its $S_n$-orbit, quotiented by the free label action of $S_n$, is
the one-object groupoid with automorphism group $B_n$. After full
completion its group is $\widehat B_n$, by (C.31.1). Since $\Phi$
fixes objects and commutes with this symmetric-group action, it
induces an actual continuous endomorphism
\[
                \phi_n:\widehat B_n\longrightarrow\widehat B_n
\]
whose map on $S_n$ is the identity. Taking the automorphism group of
just one labelled object would instead give the pure braid group;
the orbit quotient is essential here.

We next verify the exact full-twist formula without assuming that
$\phi_n$ is onto. Let $\vartheta$ be the one-strand ribbon twist.
The profinite ribbon-extension isomorphism
\cite[Proposition 8.1]{CEBHR2019} extends $\Phi$ uniquely to an
object-fixing endomorphism of $\widehat{\mathrm{PaRB}}$, preserving
the untwisted braids. The two-strand balancing calculation in
\cite[Proposition 7.7]{CEBHR2019}, valid over
$\widehat{\mathbb Z}$, gives
$\Phi(\vartheta)=\vartheta^\lambda$. More explicitly, if its
twist exponent is $t$, that calculation gives $2t=2\lambda$;
multiplication by $2$ is injective on $\widehat{\mathbb Z}$.
Put
\[
 z_n=(\sigma_1\cdots\sigma_{n-1})^n\in Z(B_n).
\]
The defining ribbon-cabling rule
\cite[Definition 6.4]{CEBHR2019} says literally
\[
 \vartheta\circ_1\operatorname{id}_b=(z_n,[1,\ldots,1]).
\]
Applying the ribbon endomorphism to the two sides gives respectively
$(z_n^\lambda,[\lambda,\ldots,\lambda])$ and
$(\phi_n(z_n),[\lambda,\ldots,\lambda])$. Consequently
\[
                         \phi_n(z_n)=z_n^\lambda.                \tag{C.31.3}
\]
This uses the operadic cabling identity; centrality of $z_n$ alone
would not justify (C.31.3) for an unproved automorphism.

Suppose now that $\lambda\in p\widehat{\mathbb Z}$ for some prime
$p$, and write $\lambda=p\nu$. In $B_p$ let
$\delta_p=\sigma_1\cdots\sigma_{p-1}$, so that
$\delta_p^p=z_p$. Centrality of $z_p$ persists in the full completion.
By (C.31.3) the element
\[
                  a=\phi_p(\delta_p)z_p^{-\nu}
\]
satisfies $a^p=1$. Its permutation is the nontrivial $p$-cycle of
$\delta_p$, since $z_p$ is pure and $\phi_p$ preserves permutations.
It therefore has exact order $p$, contradicting (C.31.2). Every
prime has been excluded, so $\lambda$ is a full unit.

Finally the full monoid is precisely the object-fixing operad
endomorphism monoid by \cite[Proposition 7.3]{CEBHR2019}. Its pairs
now have unit multiplier; C.30 makes their normalized free-group
endomorphisms automorphisms, and the closed-submonoid argument there
supplies inverses satisfying the same full equations.
\end{proof}

The use of arity $p$ for every prime is substantive. An action known
only in bounded arities would rule out only the corresponding
primes. Neither a pro-$p$ braid quotient nor a prounipotent action
supplies (C.31.1)--(C.31.2). In particular C.31 does not promote the
bare-pentagon element of A.37 to a full operad endomorphism: the full
hexagon and an actual profinite parameter still have to be supplied.

\statement{C.32}{Proved}{The relative pro-$p$ braid monoid has unit parameter}
Fix a prime $p$. Suppose that
$\mu\in1+2\mathbb Z_p$ and $f\in(F_2^{(p)})'$ satisfy the duality,
hexagon and pentagon equations in the corresponding pro-$p$ pure
groups. Then $\mu\in\mathbb Z_p^\times$. In particular the pro-$p$
Grothendieck--Teichm\"uller endomorphism monoid is its unit group.

\begin{proof}
The case $p=2$ follows already from $\mu\in1+2\mathbb Z_2$.
We give the relative braid argument, including its topology and
cohomology, for every prime. Define
\[
 B_n^{[p]}=\widehat B_n/
       \ker(\widehat P_n\longrightarrow P_n^{(p)}).
\]
The kernel being divided out is characteristic in $\widehat P_n$
and therefore normal in $\widehat B_n$. Thus
\[
 1\longrightarrow P_n^{(p)}\longrightarrow B_n^{[p]}
       \longrightarrow S_n\longrightarrow1                    \tag{C.32.1}
\]
is exact. This is a relative completion retaining $S_n$, not the
maximal pro-$p$ quotient of $B_n$.

The Fadell--Neuwirth action of $P_{n-1}$ on its free kernel
$F_{n-1}$ fixes every meridian in abelianization. This is the
almost-direct-product structure recorded explicitly in
\cite[Example 1.1]{CECohen2010}. Given a characteristic finite
$p$-group quotient $F_{n-1}/W$, the image of $P_{n-1}$ in its
automorphism group acts trivially on the Frattini quotient and
is therefore a finite $p$-group. The resulting finite semidirect
product is a finite $p$-group quotient of $P_n$ detecting
$F_{n-1}/W$. Characteristic finite $p$-quotients are cofinal for a
finitely generated free group: intersect the finitely many
automorphic conjugates of a normal subgroup of a given finite
$p$-power index. Consequently the free kernel has its full
induced pro-$p$ topology and
\[
 1\longrightarrow F_{n-1}^{(p)}\longrightarrow P_n^{(p)}
       \longrightarrow P_{n-1}^{(p)}\longrightarrow1             \tag{C.32.2}
\]
is exact and split.

Here ``$p$-good'' will mean that cohomology comparison is an
isomorphism for every finite $p$-primary module whose action extends
continuously to the pro-$p$ completion. Finite-rank free groups have
this property: their degree-zero and degree-one comparisons follow
from the free universal property, and both higher cohomologies
vanish. In (C.32.2), free-kernel cohomology with a finite such module
is finite and vanishes above degree one. Its continuous counterpart
is a finite continuous module for $P_{n-1}^{(p)}$. The
Hochschild--Serre spectral-sequence comparison therefore proves
inductively that every $P_n$ is $p$-good in precisely this sense.

Let $S$ be a Sylow $p$-subgroup of $S_n$, and let $G$ be its
preimage in the discrete group $B_n$. The normal subgroup $P_n$
has $p$-power index in $G$. If $W\triangleleft P_n$ has finite
$p$-power index, its finitely many $G/P_n$ conjugates have an
intersection normal in $G$, still of $p$-power index. Hence the
pro-$p$ topology of $G$ induces the full pro-$p$ topology on $P_n$,
and
\[
 1\longrightarrow P_n^{(p)}\longrightarrow G^{(p)}
       \longrightarrow S\longrightarrow1                       \tag{C.32.3}
\]
is exact. Comparing spectral sequences with finite quotient $S$
now proves $p$-goodness of $G$. Its configuration-space cover is
an aspherical manifold of real dimension $2n$, so
$\operatorname{cd}(G)\leq2n$. Goodness, followed by the finite-module
filtered-colimit argument of C.31, gives
\[
                    \operatorname{cd}_p(G^{(p)})\leq2n.
\]
In particular $G^{(p)}$ has no element of order $p$. The preimage
of $S$ in (C.32.1) is a pro-$p$ group and is exactly $G^{(p)}$:
the natural map is the identity on the full kernel $P_n^{(p)}$
and on the quotient $S$ in (C.32.3). Thus every Sylow $p$-subgroup
of $B_n^{[p]}$ is torsion-free.

Complete the automorphism groups of the parenthesized-braid
groupoids pro-$p$, retaining their labelled object sets. Operadic
composition extends because pro-$p$ completion commutes with finite
products. The groupoid presentation in
\cite[Lemma 7.1 and Proposition 7.3]{CEBHR2019} gives an actual
object-fixing endomorphism from our pair: the elementary braiding
goes to $\beta(\beta^2)^{(\mu-1)/2}$ and the associativity arrow
has parameter $f$. Here $\beta^2$ denotes the pure loop obtained
by composing $\beta$ with its label-swapped mate; the expression
is read in $B_2^{[p]}$, or in the corresponding morphism torsor.
Its defining relations are exactly the two
hexagons and pentagon; the homomorphisms on the pure groups factor
through their pro-$p$ completions since the target pure groups
are pro-$p$. Taking the free label-orbit quotient now gives
endomorphisms $\phi_n$ of $B_n^{[p]}$ preserving $S_n$.
The same ribbon presentation and cabling computation as in C.31
apply over $\mathbb Z_p$ and give
\[
                             \phi_n(z_n)=z_n^\mu.                \tag{C.32.4}
\]
Only powers of pure braids have $\mathbb Z_p$ exponents; the formula
for the elementary braiding above explicitly retains its odd
permutation class, even when $p$ is odd and $\mu=0$.

If $\mu=p\nu$, take $n=p$ and
$a=\phi_p(\delta_p)z_p^{-\nu}$, where
$\delta_p=\sigma_1\cdots\sigma_{p-1}$ and $\delta_p^p=z_p$.
Equation (C.32.4) gives $a^p=1$, while its permutation is a
nontrivial $p$-cycle. Choosing $S$ to contain that cycle places
$a$ in the torsion-free group $G^{(p)}$, a contradiction.
Thus $\mu$ is a unit. The free pro-$p$ Frattini criterion makes
the normalized free-group endomorphism an automorphism, and
closedness of the solution monoid gives its inverse within the
same solution set, as in C.30.
\end{proof}

\statement{C.33}{Proved}{An integral pro-$p$ pentagon has a square-unit parameter}
Let $f\in F_2^{(p)}$ satisfy the pentagon in $P_4^{(p)}$, and use
the convention $[x,y]=xyx^{-1}y^{-1}$ for its coefficient $c_2(f)$.
Then
\[
               D=1+24c_2(f)\in(\mathbb Z_p^\times)^2.            \tag{C.33.1}
\]
For either root $\mu\in\mathbb Z_p^\times$, the pair $(\mu,f)$
satisfies the actual pro-$p$ hexagons and belongs to the pro-$p$
Grothendieck--Teichm\"uller group. In particular the canonical
full pentagon word in A.37 has a peripheral conjugator with that
prescribed word at every maximal pro-$p$ quotient.

\begin{proof}
Furusho's theorem makes $f$ a commutator word and supplies duality
and the hexagon in the free prounipotent group over
$E=\mathbb Q_p(\mu)$, where $\mu^2=D$
\cite[Theorem 2 and Corollary 12]{Furusho2010}. For $p=2$,
$D\equiv1\pmod8$, and for $p=3$, $D\equiv1\pmod3$; these are
already square units. Suppose henceforth that $p\geq5$ and
$D\ne0$. Both $\mu$ and $m=(\mu-1)/2$ belong to the valuation
ring $\mathcal O_E$. We shall prove that $D$ is a nonzero square
modulo $p$ using integral metabelian coefficients. No integrality
of logarithmic or gamma-series coefficients is assumed.

Put $\Lambda=\mathbb Z_p[[A,B]]$, $a=1+A$, $b=1+B$. In the
pro-$p$ semidirect product of the additive module
$\Lambda e_1\oplus\Lambda e_2$ by $\mathbb Z_p^2$, write
\[
 (d,v)(e,w)=(de,v+dw),\qquad
 x=(a,e_1),\quad y=(b,e_2).
\]
The diagonal elements are $d=a^s b^t$ with $s,t\in\mathbb Z_p$.
The commutator is
\[
 [x,y]=(1,w_0),\qquad
 w_0=(1-b)e_1+(a-1)e_2.
\]
There is a unique series $F(a,b)\in\Lambda$ such that the image
of $f$ is $(1,F(a,b)w_0)$. Indeed the closed commutator subgroup
is normally generated by $[x,y]$, conjugation multiplies its
translation by $a^s b^t$, and the module $\Lambda w_0$ is compact
and therefore closed. Uniqueness follows because $\Lambda$ is a
domain and $w_0$ has a nonzero coordinate. This constructs the
needed coefficient; no assertion identifying arbitrary closed
subgroups or the whole derived module is needed. Duality gives
\[
                              F(a,b)=F(b,a).                    \tag{C.33.2}
\]
For arbitrary affine matrices $(d,v)$ and $(e,w)$ the evaluated
commutator has translation $(1-e)v+(d-1)w$, and the evaluated
word $f$ has this translation multiplied by $F(d,e)$. The formula
holds first for discrete words and then by continuity. Its image
under the canonical embedding in the prounipotent group is the
same Magnus expansion, coefficient by coefficient.

Set now $a=1+T$ and evaluate the prounipotent identities at
\[
 x=(a,T),\quad y=(a^{-1},0),\quad
 xy=(1,T),\quad z=(xy)^{-1}=(1,-T).                             \tag{C.33.3}
\]
These matrices lie in $1+T M_2(E[[T]])$, so logarithms and every
prounipotent substitution are defined. Write
\[
 t(a)=(1-a^{-1})F(a,a^{-1}),\qquad V(a)=(a-1)F(1,a).
\]
The three word evaluations and the powers in the hexagon are
\[
 \begin{array}{lll}
 f(x,y)=(1,Tt(a)),& f(z,x)=(1,TV(a)),&
 f(y,z)=(1,-TV(a^{-1})),\\[2pt]
 x^m=\bigl(a^m,T(a^m-1)/(a-1)\bigr),&
 z^m=(1,-mT),&y^m=(a^{-m},0).
 \end{array}                                                    \tag{C.33.4}
\]
The last word evaluation uses (C.33.2). Multiply the six factors
in the hexagon
\[
                    f(x,y)x^m f(z,x)z^m f(y,z)y^m=1
\]
and divide the resulting translation by $Ta^m$. This gives
\[
 a^{-m}t(a)+\frac{1-a^{-m}}{a-1}
                  +V(a)-V(a^{-1})-m=0.                         \tag{C.33.5}
\]
Symmetry gives $t(a^{-1})=-a t(a)$. Substituting $a^{-1}$ for
$a$ in (C.33.5) therefore gives
\[
 -a^{m+1}t(a)+\frac{a(a^m-1)}{a-1}
                  +V(a^{-1})-V(a)-m=0.                         \tag{C.33.6}
\]
Add these identities and put $d=a^{-m}-a^{m+1}$. Since
$1-a^{-m}+a^{m+1}-a=1-a-d$, their sum is
\[
                         d\bigl(t(a)-T^{-1}\bigr)=\mu.
\]
Here $d\ne0$ because $\mu\ne0$, so division in $E((T))$ is
legitimate. It yields the exact integral identity
\[
 \begin{split}
 R(T)&:=1-T(1-a^{-1})F(a,a^{-1})\\
     &=\frac{\mu a^{(\mu-1)/2}T}{a^\mu-1}
             \ \in\ 1+T^2\mathbb Z_p[[T]].
 \end{split}                                                    \tag{C.33.7}
\]

Let $H_\mu(T)=(a^\mu-1)/(\mu T)$, whose constant term is one.
Equation (C.33.7) says $R(T)H_\mu(T)=a^m$. For
$0\leq j\leq p-2$ its coefficients
\[
 [T^j]H_\mu(T)=\frac{(\mu-1)\cdots(\mu-j)}{(j+1)!}
\]
are in $\mathcal O_E$. Also
$[T^{p-1}]a^m=\binom{m}{p-1}$ is integral, since $(p-1)!$ is
a $p$-adic unit. Comparing the coefficient of $T^{p-1}$ in the
product, using $R(T)\in\mathbb Z_p[[T]]$, forces
\[
     [T^{p-1}]H_\mu(T)
       =\frac{\prod_{j=1}^{p-1}(\mu-j)}{p!}\in\mathcal O_E.
\]
For $a^m$ only degrees at most $p-1$ were used, whose factorial
denominators are prime to $p$. Integrality of that series in every
degree was not assumed. Thus
\[
 \mu^{p-1}-1=D^{(p-1)/2}-1\in p\mathcal O_E\cap\mathbb Q_p
                                      =p\mathbb Z_p,
\]
because $\prod_{j=1}^{p-1}(X-j)\equiv X^{p-1}-1\pmod p$.
Euler's criterion and Hensel's lemma now prove (C.33.1).

It remains to exclude $D=0$, rather than taking a limit of
(C.33.7) with the word $f$ held fixed. In that case Furusho's
root is the actual integral parameter $\mu=0$. Since $p\geq5$,
also $m=-1/2\in\mathbb Z_p$. The prounipotent hexagons therefore
descend to the free pro-$p$ group under its injective canonical
embedding; the given pentagon remains a pro-$p$ identity.
C.32 then excludes this nonunit parameter. The same injective
comparison for the square-unit roots already obtained shows that
the hexagons hold in the actual pro-$p$ group. Frattini and the
closed solution-monoid argument give the final unit assertion.
For the canonical word of A.37, its intertwining identity is
already a full identity and hence descends to every pro-$p$
quotient, giving the stated conjugators.
\end{proof}

\paragraph{Literature comparison and remaining full equation.}
Furusho explicitly allowed a quadratic extension for the root and
observed that it might fail to be a $p$-adic unit
\cite[paragraph following Corollary 13]{Furusho2010}.
C.32--C.33 address the integral pro-$p$ word using relative braid
torsion and the displayed Magnus coefficient. The exact full and
relative completions and all coefficient comparisons have been
given, and no literature priority is asserted. These local roots
assemble to a full unit in $\widehat{\mathbb Z}$, but the resulting
hexagon is presently established only in every maximal pro-$p$
quotient. A.39's mixed-prime example explains why this alone does
not imply the full profinite third-cusp equation or marked
involution conjugacy.

\statement{C.34}{Refuted}{A specified full square root need not satisfy the hexagon}
The literal prescribed-root assertion in Furusho's Question 14
\cite{Furusho2010} is false: a full profinite pentagon word $f$ and
a specified unit $\mu$ satisfying $\mu^2=1+24c_2(f)$ need not
satisfy the full hexagon. This does not refute the existential
assertion that some choice of such a root satisfies it.

\begin{proof}
Take $f=1$. In $\widehat{\mathbb Z}=\prod_p\mathbb Z_p$ define
$\mu_2=-1$ and $\mu_p=1$ for every odd prime $p$. Then $\mu$ is a
unit, $\mu^2=1=1+24c_2(f)$, and $f$ satisfies the full pentagon
and duality. Moreover
\[
 \mu\equiv7\pmod{24},\qquad m=(\mu-1)/2\equiv3\pmod{12}.
\]
Use the quotient of the full free group into $S_4$ given, with
right-to-left permutation products, by
\[
                     x=(234),\qquad y=(1234).
\]
The orders of $x,y,xy$ are respectively $3,4,4$. Since
$z=(xy)^{-1}$, the claimed hexagon with $f=1$ has image
\[
             x^m z^m y^m=1\cdot(xy)\cdot y^{-1}=x\ne1.
\]
This is the same finite permutation pair used in A.39 to show that
all maximal pro-$p$ peripheral tests are insufficient. Here it
directly evaluates the hexagon of the prescribed root.
For the same word $f=1$, either of the two global roots $1,-1$
does satisfy that equation. Thus the example concerns the choice
of root, not existence of a valid choice.
\end{proof}

The exact affine multiplications used in C.33 were independently
checked symbolically in
\path{code/check_metabelian_hexagon_identity.py}; C.34's actual
hexagon value is reproduced in
\path{work_CE/code/check_prescribed_root_hexagon.py}.

\statement{C.35}{Proved}{The canonical word has full duality and passes every tetrahedral generation test}
Let \(\alpha=(-1,f)\in\widehat{GT}_0\) be a nonidentity
involution, and let \(u\in F'\) be its unique R1 parameter,
\(u=f\alpha(u)\). Then
\[
                         u(y,x)=u(x,y)^{-1}.
 \tag{C.35.1}
\]
For every continuous epimorphism \(q:F\twoheadrightarrow A_4\),
\[
           \big\langle q(x),q(u)^{-1}q(y)q(u)\big\rangle=A_4.
 \tag{C.35.2}
\]
Thus the full canonical substitution passes every \(A_4\) generation
test. No pentagon hypothesis is needed for this assertion, and no
full generation or third-cusp assertion is inferred from it.

\begin{proof}
\emph{Canonical duality from full uniqueness.}\par
Let \(\iota\) interchange \(x,y\). The two-cycle relation for
\(f\) gives \(\iota(f)=f^{-1}\), and direct evaluation on the
free generators gives
\[
                   \iota\alpha=\operatorname{Inn}(f)\alpha\iota.
\]
Applying \(\iota\) to \(u=f\alpha(u)\) therefore yields
\(\iota(u)=\alpha(\iota(u))f^{-1}\). Consequently
\((\iota(u))^{-1}=f\alpha((\iota(u))^{-1})\).
Both candidate solutions belong to \(F'\). The full uniqueness in
R1 proves (C.35.1). This argument does not pass uniqueness to a
finite quotient.

\emph{The binary tetrahedral calculation.}\par
All following matrices have entries in \(\mathbb F_3\). Put
\[
\begin{gathered}
 L=\operatorname{SL}_2(\mathbb F_3),\qquad
 Q=L'=\{\pm1,\pm I,\pm J,\pm K\},\\
 I=\begin{pmatrix}0&1\\2&0\end{pmatrix},\quad
 J=\begin{pmatrix}1&1\\1&2\end{pmatrix},\quad
 K=IJ=\begin{pmatrix}1&2\\2&2\end{pmatrix},\\
 A=\begin{pmatrix}1&1\\0&1\end{pmatrix},\quad
 B=\begin{pmatrix}0&1\\2&1\end{pmatrix},\quad
 D=\begin{pmatrix}0&1\\2&2\end{pmatrix}.
\end{gathered}
 \tag{C.35.3}
\]
Here \(I^2=J^2=K^2=-1\), and \(L/\{\pm1\}\simeq A_4\).
The subgroup \(Q\) consists of \(\pm1\) and all six trace-zero
matrices, so it is normal and \(L/Q\) has order three. Thus
\(L'\subseteq Q\); conversely \([A,I]=-K\) and its
\(A\)-conjugates generate \(Q\), proving \(L'=Q\).
The four pairs
\[
              (A,I),\quad(A,B),\quad(A,D),\quad(I,A)
 \tag{C.35.4}
\]
represent the four automorphism orbits of generating pairs in
\(A_4\). Their projective passports are respectively
\((3,2,3),(3,3,2),(3,3,3),(2,3,3)\).
Indeed the mixed-order pairs give two orbits of size twenty-four.
For two elements of order three, belonging to the same or opposite
nontrivial class in \(A_4^{\rm ab}=C_3\) gives the other two
orbits, again of size twenty-four. Conjugation by
\(\operatorname{GL}_2(\mathbb F_3)\) induces every automorphism
of \(A_4\).

For any commutator word \(w\), denote its values at (C.35.4) by
\(W_0,W_1,W_2,W_3\in Q\). Multiplying either input by a
central element \(-1\) leaves \(w\)'s value unchanged, since
its two abelian exponents are zero. This remains true for a full
profinite word: the central correction factors through its
continuous abelianization.

Suppose \(w\) satisfies duality. Then
\[
 W_3=W_0^{-1},\qquad
 W_1\in\{\pm1,\pm K\},\qquad
 W_2\in\{\pm1,\pm I,\pm K\}.
 \tag{C.35.5}
\]
To check the second assertion, conjugation by
\[
                         S=\begin{pmatrix}1&0\\2&2\end{pmatrix}
\]
sends \((A,B)\) to \((-B,-A)\) and acts on
\((I,J,K)\) as \((-J,-I,-K)\). Thus duality requires
\(W_1\operatorname{Ad}(S)(W_1)=1\); its solutions in
\(Q\) are exactly \(\pm1,\pm K\).
For the third assertion, conjugation by \(J\) interchanges
\(A,D\) and acts on \((I,J,K)\) as
\((-I,J,-K)\). The analogous equation excludes precisely
\(\pm J\). These checks involve all eight elements of \(Q\).

Apply this to both \(f\) and \(u\), writing their values as
\(F_j\) and \(U_j\). The negative hexagon for \(f\) is
\[
       f(a,b)a^{-1}f(z,a)z^{-1}f(b,z)b^{-1}=1,
       \qquad z=(ab)^{-1}.
 \tag{C.35.6}
\]
At \((a,b)=(A,B)\), one has \(Z=(AB)^{-1}=J\),
\(AIA^{-1}=-J\), and
\(S(A,I)S^{-1}=(-B,-J)\). Therefore
\[
          f(Z,A)=AF_3A^{-1},\qquad f(B,Z)=SF_0S^{-1}.
\]
Since \(A^{-1}Z^{-1}=B\), equation (C.35.6) becomes
\[
                        F_1F_3\rho(F_0)=1,
 \qquad \rho=\operatorname{Ad}(R),\quad
 R=1+I=\begin{pmatrix}1&1\\2&1\end{pmatrix},
 \tag{C.35.7}
\]
where \(BS=-R\). The automorphism \(\rho\) fixes \(I\)
and sends \(J\) to \(K\), and \(K\) to \(-J\).
Using \(F_3=F_0^{-1}\), we obtain
\[
                    F_1=\rho(F_0)^{-1}F_0\in\{1,I\}.
\]
For \(F_0\in\{\pm1,\pm I\}\) the value is one; for
\(F_0\in\{\pm J,\pm K\}\) it is \(I\).
Its intersection with the second set in (C.35.5) is just one.
Consequently
\[
                              f(A,B)=F_1=1.
 \tag{C.35.8}
\]

Now evaluate the full coboundary equation at \((A,B)\).
By (C.35.8) it says
\[
                          U_1=u(A^{-1},B^{-1}).
\]
Conjugation by
\[
                         T=\begin{pmatrix}1&2\\0&2\end{pmatrix}
\]
sends \((A,B)\) to \((A^{-1},B^{-1})\) and acts on
\((I,J,K)\) as \((J,I,-K)\). Hence \(U_1\) is fixed
by this automorphism. In the set \(\{\pm1,\pm K\}\)
allowed by duality, the fixed elements are precisely \(\pm1\).
We have proved
\[
                              U_1\in\{\pm1\}.
 \tag{C.35.9}
\]
Thus the projective \((3,3,2)\) generating pair is unchanged
by the canonical substitution.

For the \((3,3,3)\) pair, the only projective quaternion
conjugator that collapses \((\bar A,\bar D)\) into one
subgroup of order three is \(\bar J\): in fact
\(J^{-1}DJ=A\). For conjugators \(1,I,K\), the second
projective element belongs to a different subgroup of order
three, so the pair generates \(A_4\). This last assertion
can also be read directly from the four projective values in
(C.35.3). Duality has already excluded \(U_2\in\{\pm J\}\).

The two mixed-order cases are immediate. Conjugation by an
element of the projective \(V_4\) fixes every element of
\(V_4\); and any element of order two together with an
element of order three generates \(A_4\). The four orbits
(C.35.4) exhaust every epimorphism \(F\to A_4\).
Lifting the two chosen generators arbitrarily to \(L\), and
using equivariance and invariance under central changes of the
inputs, transports the calculation to each such epimorphism.
This proves (C.35.2).
\end{proof}

The literal matrix identities and the four projective generation
checks are independently reproduced in
\path{work_CE/code/check_C35_matrices.py}.
The longer files \path{work_CE/code/check_a4_pentagon.py} and
\path{work_CE/code/check_sl23_pentagon.py} record the discovery
calculation and its complete finite consistency checks. The
necessary-condition proof above uses the given full words and
only their displayed identities; it does not assume that a
finite word-evaluation table lifts to a full solution.

\statement{C.36}{Proved}{Tetrahedral third-cusp tests and dyadic generation over the full modulo-three base}
Retain the hypotheses and canonical word \(u\) of C.35. For every
\(\mu\in\widehat{\mathbb Z}^{\times}\) and every continuous
homomorphism \(q:F\to A_4\), one has
\[
 q\bigl(x^\mu u^{-1}y^\mu u\bigr)
                  \sim_{A_4}q\bigl((xy)^\mu\bigr).
 \tag{C.36.1}
\]
Put
\[
 U_3=\ker\bigl(F\to C_3^2\bigr)=\overline{F'F^3},\qquad
                      G_{3,2}=F/O^2(U_3).
 \tag{C.36.2}
\]
For every such \(\mu\), the normalized substitution
\[
                   x\longmapsto x^\mu,\qquad
                   y\longmapsto u^{-1}y^\mu u
 \tag{C.36.3}
\]
induces an automorphism of the actual relative pro-two group
\(G_{3,2}\), and intertwines the induced actions of the given
\(\alpha\) and standard \(c\). In particular this proves
the canonical generation condition at the specified
\(U=F'F^3\) relative level. The full relative third-cusp
condition is not asserted.

\begin{proof}
\emph{The finite third cusp.}\par
Every proper subgroup of \(A_4\) is abelian. If \(q\) is
not onto, \(q(u)=1\) and (C.36.1) is equality.
Suppose \(q\) is onto. By C.35, and since a unit power
generates the same finite cyclic subgroup, the two substituted
projective generators still generate \(A_4\).

If \(q(xy)\) has order three, its nontrivial image in
\(A_4^{\rm ab}=C_3\) determines its entire conjugacy class.
The substituted product has the same \(\mu\)-powered
abelian class as \(q((xy)^\mu)\), so the two are conjugate.
If \(q(xy)\) has order two, then \(q(x),q(y)\) have
order three and opposite nonzero abelian classes. The new
product lies in \(V_4\), and it is not one because the new
pair generates \(A_4\). All three nonidentity elements of
\(V_4\) are conjugate in \(A_4\), proving (C.36.1).
These cases exhaust the generating pairs.

\emph{The exact relative Frattini module.}\par
The full open subgroup \(U_3\) has index nine and is free
profinite of rank ten. Thus \(N=U_3/O^2(U_3)\) is free
pro-two of rank ten, and
\[
                    1\longrightarrow N\longrightarrow G_{3,2}
                       \longrightarrow Q=C_3^2\longrightarrow1
\]
is the defining relative extension. Let
\(M=N/\Phi(N)=H_1(U_3,\mathbb F_2)\).
The cellular chain sequence of the nine-sheeted Cayley graph
of the two-petalled rose is
\[
 0\longrightarrow M\longrightarrow\mathbb F_2[Q]^2
   \longrightarrow\mathbb F_2[Q]\longrightarrow\mathbb F_2
   \longrightarrow0.
\]
Since nine is prime to two, this sequence splits as modules
where appropriate, and gives
\[
 M\simeq\mathbb F_2\oplus\mathbb F_2[Q]
       \simeq\mathbb F_2^2\oplus W_1\oplus W_2\oplus W_3\oplus W_4.
 \tag{C.36.4}
\]
Here the \(W_i\) are the four pairwise nonisomorphic
nontrivial simple \(\mathbb F_2[Q]\)-modules, each of
dimension two. Over \(\mathbb F_4\) the eight nontrivial
characters pair with their Frobenius conjugates; each pair
has the same one-dimensional kernel in \(Q\), giving
precisely these four modules.

First take \(\mu=1\) and write \(\psi\) for the
canonical substitution. It preserves \(U_3\) and its
fully invariant pro-two residual, so it induces an
endomorphism of \(G_{3,2}\). Its action on \(Q\) is
identity, and hence its induced linear map on \(M\) is
\(Q\)-equivariant: changing a lift of an element of
\(Q\) by an element of \(N\) has no effect on
\(N/\Phi(N)\). The invariant summand in (C.36.4) is
identified with \(H_1(F,\mathbb F_2)\) by transfer, or
by the coinvariant sequence for the odd-order quotient
\(Q\). On this summand \(\psi\) is identity.

For each \(W_i\), the finite group
\(G_{3,2}/\Phi(N)\) has a quotient
\[
                        W_i\rtimes C_3\simeq A_4.
 \tag{C.36.5}
\]
Indeed its extension by \(Q\) splits by the coprime
complement theorem; quotient by the other summands and then
by the kernel of the action of \(Q\) on \(W_i\).
If \(\psi\) were zero on \(W_i\), its composite with
(C.36.5) would kill \(N\), and could not be onto \(A_4\).
That contradicts C.35. By simplicity, its restriction to
each \(W_i\) is therefore invertible. It follows that
\(\psi\) is onto \(M\), so Frattini gives
\(\psi(N)=N\). It is also onto \(Q\), hence onto
\(G_{3,2}\). This finitely generated profinite group is
Hopfian, so the induced endomorphism is an automorphism.

For an arbitrary full unit \(\mu\), (C.36.3) still
preserves \(U_3\) and \(O^2(U_3)\). Its two images
generate exactly the same closed subgroup of \(G_{3,2}\)
as \(x,u^{-1}yu\), because unit powering preserves the
whole closed procyclic subgroup. Thus it too is onto and
is an automorphism. Finally \(u=f\alpha(u)\) gives
\(\alpha(\mu,u)=(\mu,u)c\) on the two free generators,
so the stated intertwining holds in every indicated quotient.
Nothing in this Frattini argument identifies the relative
conjugacy class of the third cusp.
\end{proof}
